% ============================================================
%  Trabajo final completo (las 4 partes). Cada parte también se compila sola:
%  RLM.tex, Logistica.tex, CP.tex y Conglomerados.tex (mismo preámbulo y secciones/)
% ============================================================
\documentclass[11pt,a4paper]{article}
% ============================================================
%  Basado en la plantilla de las Prácticas 8-11 (mismos colores/estilos)
% ============================================================

% ---------- Idioma y codificación ----------
\usepackage[utf8]{inputenc}
\usepackage[T1]{fontenc}
\usepackage[spanish,es-noshorthands,es-tabla]{babel}

% ---------- Tipografía y layout ----------
\usepackage{lmodern}
\usepackage[sfdefault,light]{roboto}
\usepackage{geometry}
\geometry{a4paper, top=3.1cm, bottom=2.8cm, left=2.4cm, right=2.4cm, headheight=1.6cm}
\usepackage{ragged2e}
\usepackage{enumitem}
\usepackage{graphicx}
\usepackage{float}
\usepackage{booktabs}
\usepackage{array}
\usepackage[most]{tcolorbox}
\usepackage{tikz}
\usetikzlibrary{shadows,calc}
\usepackage{titlesec}
\usepackage{fancyhdr}
\usepackage{fontawesome5}
\usepackage[colorlinks=true,breaklinks=true]{hyperref}
\usepackage{caption}
\usepackage{microtype}
\usepackage{xurl}
\usepackage{listings}
\usepackage{amsmath,amssymb}
\usepackage{needspace}
\usepackage{icomma}

% ---------- Paleta de colores (tema "coolnight" de tu terminal Kitty) ----------
\definecolor{cNavy}{HTML}{011423}
\definecolor{cTealDeep}{HTML}{064B5A}
\definecolor{cTeal}{HTML}{097289}
\definecolor{cBlue}{HTML}{0FC5ED}
\definecolor{cCyan}{HTML}{0FC5ED}
\definecolor{cMint}{HTML}{44FFB1}
\definecolor{cPurple}{HTML}{A277FF}
\definecolor{cRed}{HTML}{E52E2E}
\definecolor{cBg}{HTML}{EAF6F8}
\definecolor{cText}{HTML}{01111E}
\definecolor{cOutBg}{HTML}{F4F4F4}

\hypersetup{
  linkcolor=cTealDeep,
  citecolor=cTealDeep,
  urlcolor=cBlue,
  pdftitle={Trabajo final: regresión lineal y logística, componentes principales y conglomerados},
  pdfauthor={Miguel Deza}
}

\color{cText}
\emergencystretch=2.5em

% ---------- Encabezado / pie de página ----------
\pagestyle{fancy}
\fancyhf{}
\renewcommand{\headrulewidth}{0pt}
\renewcommand{\footrulewidth}{0pt}
\fancyhead[L]{\small\color{cTealDeep}\textbf{\faGraduationCap\ Maestría en Ciencia de Datos}}
\fancyhead[R]{\small\color{cTeal}Estadística en Ciencia de Datos -- Trabajo final}
\fancyfoot[C]{%
  \begin{tikzpicture}[remember picture]
    \node[circle, fill=cTealDeep, minimum size=1.7em, inner sep=0pt] (pg) {\color{white}\small\thepage};
  \end{tikzpicture}%
}
\renewcommand{\headrule}{{\color{cCyan}\hrule height 1.4pt}\vspace{1pt}{\color{cBlue}\hrule height 0.5pt}}

% ---------- Estilo de secciones ----------
\titleformat{\section}
  {\normalfont\Large\bfseries\color{cTealDeep}}
  {\tikz[baseline=-0.6ex]{\node[fill=cCyan, text=white, rounded corners=1.5pt, inner sep=3pt, font=\bfseries]{\thesection};}}
  {0.6em}{}
  [{\vspace{-0.3em}\color{cCyan}\titlerule[1.2pt]\vspace{0.3em}}]

\titleformat{\subsection}
  {\normalfont\large\bfseries\color{cTeal}}
  {\color{cTeal}\thesubsection}{0.6em}{}

\titlespacing*{\section}{0pt}{1.4em}{0.8em}
\titlespacing*{\subsection}{0pt}{1.1em}{0.5em}

\captionsetup{labelfont={bf,color=cTealDeep}, font=small, textfont=it}

% ---------- Bloques de código R y salida de consola ----------
\lstdefinestyle{styleR}{
  language=R,
  backgroundcolor=\color{cBg},
  basicstyle=\ttfamily\footnotesize\color{cText},
  keywordstyle=\color{cTealDeep}\bfseries,
  commentstyle=\color{gray}\itshape,
  stringstyle=\color{cTeal},
  numbers=none,
  breaklines=true,
  frame=l,
  framerule=2.2pt,
  rulecolor=\color{cCyan},
  xleftmargin=8pt,
  framexleftmargin=6pt,
  aboveskip=6pt, belowskip=6pt,
  showstringspaces=false,
  columns=fullflexible, keepspaces=true
}

\lstdefinestyle{styleOut}{
  backgroundcolor=\color{cOutBg},
  basicstyle=\ttfamily\scriptsize\color{cText},
  breaklines=true,
  frame=l,
  framerule=2.2pt,
  rulecolor=\color{cTealDeep},
  xleftmargin=8pt,
  framexleftmargin=6pt,
  aboveskip=2pt, belowskip=6pt,
  showstringspaces=false,
  columns=fullflexible, keepspaces=true
}

% ---------- Preguntas ----------
\setlist{itemsep=1pt, topsep=3pt}
\graphicspath{{img/}}
\setlength{\parskip}{3pt}

\newtcolorbox{enunciado}[1][]{enhanced, breakable, colback=cBg, colframe=cTealDeep,
  boxrule=0pt, leftrule=3pt, arc=2pt, left=8pt, right=8pt, top=5pt, bottom=5pt,
  fonttitle=\bfseries, coltitle=white, colbacktitle=cTealDeep,
  title={Pregunta #1}}
\newtcolorbox{resultado}{enhanced, breakable, colback=cMint!12, colframe=cTealDeep,
  boxrule=0.5pt, arc=2pt, left=8pt, right=8pt, top=4pt, bottom=4pt}

% \pregunta{etiqueta}{texto del índice}{enunciado}
\newcommand{\pregunta}[3]{%
  \Needspace{12\baselineskip}\medskip
  \phantomsection
  \addcontentsline{toc}{subsection}{#1 \ #2}%
  \begin{enunciado}[#1]#3\end{enunciado}\noindent}

% portada de sección: página completa con fondo oscuro
\newcommand{\seccion}[4]{%
  \clearpage\newgeometry{margin=0cm}
  \refstepcounter{section}\phantomsection
  \addcontentsline{toc}{section}{\protect\numberline{\thesection}\protect\faIcon{#1}\ #2}%
  \thispagestyle{empty}\noindent
  \begin{tikzpicture}[remember picture, overlay]
    \fill[cNavy] (current page.south west) rectangle (current page.north east);
    \fill[cTealDeep, opacity=0.55] (current page.north west) ++(1.5,-1) circle (4.2);
    \fill[cCyan, opacity=0.14] (current page.north west) ++(4.2,-0.4) circle (2.2);
    \fill[cPurple, opacity=0.26] (current page.south east) ++(-1.2,1.4) circle (3.6);
    \fill[cMint, opacity=0.13] (current page.south east) ++(-3.6,0.3) circle (2.0);
    \draw[cCyan, opacity=0.25, line width=0.8pt] (current page.center) ++(5.2,4.6) circle (1.1);
    \node[anchor=center] at ($(current page.center)+(0,5.2)$)
      {\color{cTealDeep!80!cCyan}\bfseries\fontsize{150}{150}\selectfont \thesection};
    \node[circle, fill=cTealDeep, draw=cCyan, line width=1.3pt, minimum size=2.4cm] (ic)
      at ($(current page.center)+(0,1.6)$) {};
    \node at (ic) {\color{cMint}\fontsize{30}{30}\selectfont\faIcon{#1}};
    \node[anchor=center] at ($(current page.center)+(0,-0.4)$)
      {\color{cCyan}\bfseries\fontsize{11}{13}\selectfont S E C C I Ó N};
    \node[anchor=north, text width=14cm, align=center] (tit) at ($(current page.center)+(0,-0.9)$)
      {\hyphenpenalty=10000 {\color{white}\bfseries\fontsize{30}{36}\selectfont #3}};
    \draw[cCyan, line width=1.4pt] ($(tit.south)+(-3.2,-0.3)$) -- ++(6.4,0);
    \node[anchor=north, text width=14cm, align=center] at ($(tit.south)+(0,-0.6)$)
      {\color{cMint}\fontsize{13}{16}\selectfont #4};
  \end{tikzpicture}%
  \clearpage\restoregeometry}

% código y salidas: las líneas de cada bloque se leen de rangos.tex,
% generado por generar-salidas.R desde RLM.R, Logistica.R, CP.R y Conglomerados.R
\expandafter\def\csname r@rlm-datos\endcsname{28-34}
\expandafter\def\csname s@rlm-datos\endcsname{RLM.R}
\expandafter\def\csname r@rlm-ajuste\endcsname{38-41}
\expandafter\def\csname s@rlm-ajuste\endcsname{RLM.R}
\expandafter\def\csname r@rlm-r2\endcsname{45-47}
\expandafter\def\csname s@rlm-r2\endcsname{RLM.R}
\expandafter\def\csname r@rlm-supuestos\endcsname{51-53}
\expandafter\def\csname s@rlm-supuestos\endcsname{RLM.R}
\expandafter\def\csname r@rlm-diag1\endcsname{57-65}
\expandafter\def\csname s@rlm-diag1\endcsname{RLM.R}
\expandafter\def\csname r@rlm-diag2\endcsname{69-77}
\expandafter\def\csname s@rlm-diag2\endcsname{RLM.R}
\expandafter\def\csname r@rlm-vif\endcsname{81-81}
\expandafter\def\csname s@rlm-vif\endcsname{RLM.R}
\expandafter\def\csname r@rlm-backward\endcsname{85-85}
\expandafter\def\csname s@rlm-backward\endcsname{RLM.R}
\expandafter\def\csname r@rlm-final\endcsname{89-93}
\expandafter\def\csname s@rlm-final\endcsname{RLM.R}
\expandafter\def\csname r@rlm-final-sup\endcsname{97-100}
\expandafter\def\csname s@rlm-final-sup\endcsname{RLM.R}
\expandafter\def\csname r@rlm-diag-final\endcsname{104-112}
\expandafter\def\csname s@rlm-diag-final\endcsname{RLM.R}
\expandafter\def\csname r@rlm-sensibilidad\endcsname{116-121}
\expandafter\def\csname s@rlm-sensibilidad\endcsname{RLM.R}
\expandafter\def\csname r@log-datos\endcsname{28-33}
\expandafter\def\csname s@log-datos\endcsname{Logistica.R}
\expandafter\def\csname r@log-justifica\endcsname{37-46}
\expandafter\def\csname s@log-justifica\endcsname{Logistica.R}
\expandafter\def\csname r@log-ajuste\endcsname{50-53}
\expandafter\def\csname s@log-ajuste\endcsname{Logistica.R}
\expandafter\def\csname r@log-or\endcsname{57-58}
\expandafter\def\csname s@log-or\endcsname{Logistica.R}
\expandafter\def\csname r@log-lr\endcsname{62-64}
\expandafter\def\csname s@log-lr\endcsname{Logistica.R}
\expandafter\def\csname r@log-confusion\endcsname{68-74}
\expandafter\def\csname s@log-confusion\endcsname{Logistica.R}
\expandafter\def\csname r@log-roc\endcsname{78-82}
\expandafter\def\csname s@log-roc\endcsname{Logistica.R}
\expandafter\def\csname r@log-roc-fig\endcsname{86-98}
\expandafter\def\csname s@log-roc-fig\endcsname{Logistica.R}
\expandafter\def\csname r@log-cv\endcsname{102-111}
\expandafter\def\csname s@log-cv\endcsname{Logistica.R}
\expandafter\def\csname r@acp-datos\endcsname{28-34}
\expandafter\def\csname s@acp-datos\endcsname{CP.R}
\expandafter\def\csname r@acp-cor\endcsname{38-43}
\expandafter\def\csname s@acp-cor\endcsname{CP.R}
\expandafter\def\csname r@acp-corplot\endcsname{47-51}
\expandafter\def\csname s@acp-corplot\endcsname{CP.R}
\expandafter\def\csname r@acp-pca\endcsname{55-61}
\expandafter\def\csname s@acp-pca\endcsname{CP.R}
\expandafter\def\csname r@acp-scree\endcsname{65-76}
\expandafter\def\csname s@acp-scree\endcsname{CP.R}
\expandafter\def\csname r@acp-cargas\endcsname{80-81}
\expandafter\def\csname s@acp-cargas\endcsname{CP.R}
\expandafter\def\csname r@acp-cargas-fig\endcsname{85-94}
\expandafter\def\csname s@acp-cargas-fig\endcsname{CP.R}
\expandafter\def\csname r@acp-individuos\endcsname{98-114}
\expandafter\def\csname s@acp-individuos\endcsname{CP.R}
\expandafter\def\csname r@acp-biplot\endcsname{118-121}
\expandafter\def\csname s@acp-biplot\endcsname{CP.R}
\expandafter\def\csname r@acp-regresion\endcsname{125-131}
\expandafter\def\csname s@acp-regresion\endcsname{CP.R}
\expandafter\def\csname r@clu-datos\endcsname{28-37}
\expandafter\def\csname s@clu-datos\endcsname{Conglomerados.R}
\expandafter\def\csname r@clu-dist\endcsname{41-44}
\expandafter\def\csname s@clu-dist\endcsname{Conglomerados.R}
\expandafter\def\csname r@clu-dist-fig\endcsname{48-59}
\expandafter\def\csname s@clu-dist-fig\endcsname{Conglomerados.R}
\expandafter\def\csname r@clu-ward\endcsname{63-68}
\expandafter\def\csname s@clu-ward\endcsname{Conglomerados.R}
\expandafter\def\csname r@clu-dendro\endcsname{72-81}
\expandafter\def\csname s@clu-dendro\endcsname{Conglomerados.R}
\expandafter\def\csname r@clu-perfil-hc\endcsname{85-94}
\expandafter\def\csname s@clu-perfil-hc\endcsname{Conglomerados.R}
\expandafter\def\csname r@clu-kmeans\endcsname{98-109}
\expandafter\def\csname s@clu-kmeans\endcsname{Conglomerados.R}
\expandafter\def\csname r@clu-centroides\endcsname{113-121}
\expandafter\def\csname s@clu-centroides\endcsname{Conglomerados.R}
\expandafter\def\csname r@clu-compara\endcsname{125-136}
\expandafter\def\csname s@clu-compara\endcsname{Conglomerados.R}
\expandafter\def\csname r@clu-plano\endcsname{140-151}
\expandafter\def\csname s@clu-plano\endcsname{Conglomerados.R}
\expandafter\def\csname r@clu-enfermedad\endcsname{155-161}
\expandafter\def\csname s@clu-enfermedad\endcsname{Conglomerados.R}
\expandafter\def\csname r@clu-enf-fig\endcsname{165-173}
\expandafter\def\csname s@clu-enf-fig\endcsname{Conglomerados.R}

\newcommand{\codigo}[1]{%
  \edef\figopts{linerange={\csname r@#1\endcsname}}\edef\figfile{\csname s@#1\endcsname}%
  \expandafter\lstinputlisting\expandafter[\figopts, style=styleR]{\figfile}}
\newcommand{\salida}[1]{\lstinputlisting[style=styleOut]{salidas/#1.txt}}
% figura a la izquierda y su código R a la derecha
\newcommand{\figcod}[2]{%
  \par\medskip\noindent
  \begin{minipage}[c]{0.47\linewidth}\centering
    \includegraphics[width=\linewidth]{#1}\captionof{figure}{#2}
  \end{minipage}\hfill
  \begin{minipage}[c]{0.5\linewidth}
    \edef\figopts{linerange={\csname r@#1\endcsname}}\edef\figfile{\csname s@#1\endcsname}%
    \expandafter\lstinputlisting\expandafter[\figopts, style=styleR, basicstyle=\ttfamily\scriptsize\color{cText}]{\figfile}
  \end{minipage}\par\medskip}
\newcommand{\res}[1]{\begin{resultado}\textbf{\color{cTealDeep}Resultado.} #1\end{resultado}}
\newcommand{\var}[1]{\texttt{\detokenize{#1}}}

% ============================================================
% portada: \portada{icono}{título}{subtítulo}{descripción}
\newcommand{\portada}[4]{%
\begin{titlepage}
\newgeometry{margin=0cm}
\begin{tikzpicture}[remember picture, overlay]
  \fill[cNavy] (current page.south west) rectangle (current page.north east);
  \fill[cPurple, opacity=0.28] (current page.north east) ++(-1,1) circle (3.2);
  \fill[cCyan, opacity=0.22] (current page.north east) ++(-2.6,0.4) circle (2.1);
  \fill[cTealDeep, opacity=0.55] (current page.south west) ++(1,-1) circle (4.0);
  \fill[cMint, opacity=0.16] (current page.south west) ++(2.4,-0.2) circle (2.3);

  \node[circle, fill=cTealDeep, draw=cCyan, line width=1.3pt, minimum size=2.5cm] (logo) at ($(current page.north)+(0,-2.9)$) {};
  \node at (logo) {\color{cMint}\fontsize{28}{28}\selectfont\faIcon{#1}};

  \node[anchor=north, text width=14cm, align=center] at ($(current page.north)+(0,-4.7)$) {
    {\color{cCyan}\bfseries\fontsize{12}{14}\selectfont \faDatabase\ \ MAESTRÍA EN CIENCIA DE DATOS}
  };

  \node[anchor=center, text width=15cm, align=center] at ($(current page.center)+(0,1.2)$) {
    {\color{white}\bfseries\fontsize{30}{38}\selectfont #2}\\[0.35em]
    {\color{cMint}\bfseries\fontsize{17}{21}\selectfont \faIcon{#1}\ #3}
  };

  \draw[cCyan, line width=1.4pt] ($(current page.center)+(-3.6,-0.9)$) -- ++(7.2,0);

  \node[anchor=center, text width=13cm, align=center] at ($(current page.center)+(0,-1.9)$) {
    {\color{white!85!black}\fontsize{11}{14}\selectfont
     #4}
  };

  \node[anchor=south, text width=14cm, align=center] at ($(current page.south)+(0,2.6)$) {
    \begin{tabular}{r l}
      \color{cCyan}\bfseries \faUser\ Elaborado por: & \color{white} Miguel Deza \\[3pt]
      \color{cCyan}\bfseries \faGraduationCap\ Programa: & \color{white} Maestría en Ciencia de Datos -- UNSA \\[3pt]
      \color{cCyan}\bfseries \faBook\ Curso: & \color{white} Estadística en Ciencia de Datos \\[3pt]
      \color{cCyan}\bfseries \faChalkboardTeacher\ Docente: & \color{white} Carlos Antonio Abanto Valle \\[3pt]
      \color{cCyan}\bfseries \faIcon{calendar-alt}\ Fecha: & \color{white} 7 de noviembre de 2026 \\
    \end{tabular}
  };
\end{tikzpicture}
\end{titlepage}
\restoregeometry
}

\begin{document}
\portada{chart-line}{Trabajo Final\\ Estadística en Ciencia de Datos}{Modelos supervisados y no supervisados en R}{Regresión lineal múltiple y selección backward $\cdot$ Regresión logística binaria\\
     Componentes principales $\cdot$ Conglomerados Ward.D2 y $k$-means\\[2pt]
     Bases: Student Performance y Heart Disease (Cleveland), UCI}

\setcounter{tocdepth}{2}
\tableofcontents

\vspace{1.2em}
\noindent El análisis está en cuatro scripts independientes, uno por ejercicio: \var{RLM.R},
\var{Logistica.R}, \var{CP.R} y \var{Conglomerados.R} (R 4.6.1; paquetes \var{car}, \var{lmtest},
\var{pROC}, \var{cluster} y \var{corrplot}). Las bases son \var{student.csv} (395 estudiantes) y
\var{heart.csv} (297 pacientes, sin valores faltantes). En cada pregunta muestro el bloque de código, su
salida de consola y la interpretación.

\seccion{chart-line}{Regresión lineal múltiple y selección backward}{Regresión lineal\\ múltiple}{Student Performance $\cdot$ preguntas a -- e y método backward}
\section*{Datos y modelo}
\begin{align*}
\text{nota\_final}_i&=\beta_0+\beta_1\,\text{edad}_i+\beta_2\,\text{tiempo\_estudio}_i
+\beta_3\,\text{inasistencias}_i+\beta_4\,\text{reprobaciones\_previas}_i\\
&\quad+\beta_5\,\text{educacion\_madre}_i+\beta_6\,\text{educacion\_padre}_i
+\beta_7\,\text{internet\_casa}_i\\
&\quad+\beta_8\,\text{apoyo\_familiar}_i+\varepsilon_i,
\qquad \varepsilon_i\sim N(0,\sigma^2)
\end{align*}
Las variables binarias \var{internet_casa} y \var{apoyo_familiar} entran como factores con nivel de
referencia \textit{no}.
\codigo{rlm-datos}
\salida{rlm-datos}
La nota va de 0 a 20 y hay \textbf{38 notas iguales a 0}. Las revisé: los 38 tienen 0 inasistencias,
es decir, son alumnos que no rindieron la evaluación final (abandono), no notas bajas ``normales''.
Las mantengo en el ajuste porque el enunciado pide el modelo sobre la base completa, pero este
detalle explica casi todos los problemas de supuestos que aparecen después.

\pregunta{a)}{Ajuste por MCO y variables significativas}{Ajuste el modelo de regresión lineal múltiple
utilizando mínimos cuadrados ordinarios. Presente la tabla de estimación de los parámetros e identifique
las variables estadísticamente significativas al nivel $\alpha=0,05$.}
\codigo{rlm-ajuste}
\salida{rlm-ajuste}

\begin{center}\small
\begin{tabular}{lrrrrc}
\toprule
Variable & $\hat\beta$ & EE & $t$ & $p$ & ¿Signif.? \\ \midrule
Intercepto & 14,079 & 3,171 & 4,44 & $<0,001$ & sí \\
edad & $-0,306$ & 0,178 & $-1,72$ & 0,087 & no \\
tiempo\_estudio & 0,278 & 0,263 & 1,06 & 0,290 & no \\
inasistencias & 0,032 & 0,028 & 1,17 & 0,242 & no \\
reprobaciones\_previas & $-1,894$ & 0,309 & $-6,13$ & $<0,001$ & \textbf{sí} \\
educacion\_madre & 0,584 & 0,256 & 2,28 & 0,023 & \textbf{sí} \\
educacion\_padre & $-0,053$ & 0,255 & $-0,21$ & 0,834 & no \\
internet\_casa (sí) & 0,552 & 0,588 & 0,94 & 0,348 & no \\
apoyo\_familiar (sí) & $-0,984$ & 0,454 & $-2,17$ & 0,031 & \textbf{sí} \\ \bottomrule
\end{tabular}
\end{center}
El test $F$ global ($F_{8,386}=9,77$, $p=2,2\times10^{-12}$) rechaza que todos los $\beta_j$ sean cero.

\res{Al 5\,\% son significativas \textbf{reprobaciones\_previas}, \textbf{educacion\_madre} y
\textbf{apoyo\_familiar}; edad queda al borde ($p=0,087$).}

\pregunta{b)}{Interpretación de tiempo\_estudio, inasistencias y reprobaciones\_previas}{Interprete los coeficientes
asociados a las variables tiempo\_estudio, inasistencias y reprobaciones\_previas, manteniendo constantes las
demás covariables.}
\begin{itemize}[leftmargin=1.2em]
  \item[\color{cTeal}\faClock] \textbf{tiempo\_estudio} ($\hat\beta_2=0,278$): subir un nivel en la escala de
  estudio semanal (1: $<2$\,h \dots\ 4: $>10$\,h) se asocia con 0,28 puntos más en la nota final, con lo demás
  fijo. El efecto no es significativo ($p=0,29$): no puedo descartar que sea cero.
  \item[\color{cTeal}\faIcon{calendar-times}] \textbf{inasistencias} ($\hat\beta_3=0,032$): cada falta adicional
  se asocia con 0,03 puntos \emph{más}, también no significativo ($p=0,24$). El signo contraintuitivo viene de los
  38 alumnos con nota 0 y 0 faltas: la base registra 0 inasistencias para quien abandonó, y eso arrastra la
  pendiente hacia arriba.
  \item[\color{cRed}\faTimesCircle] \textbf{reprobaciones\_previas} ($\hat\beta_4=-1,894$): cada curso
  reprobado antes se asocia con 1,89 puntos menos en la nota final (casi 10\,\% de la escala), manteniendo
  constantes las demás covariables. Es el efecto más fuerte y claro del modelo ($t=-6,13$).
\end{itemize}

\pregunta{c)}{$R^2$ y $R^2$ ajustado}{Reporte el coeficiente de determinación $R^2$ y el $R^2$ ajustado.
Explique por qué el $R^2$ ajustado es más apropiado en este contexto.}
\codigo{rlm-r2}
\salida{rlm-r2}
\[
R^2=1-\frac{SCE}{SCT}=0,1684,\qquad
R^2_{aj}=1-(1-R^2)\frac{n-1}{n-p-1}=1-(0,8316)\frac{394}{386}=0,1512
\]
El modelo explica cerca del 17\,\% de la variabilidad de la nota final. El $R^2$ nunca baja al agregar
variables, aunque sean ruido; el $R^2_{aj}$ penaliza cada parámetro extra con $(n-1)/(n-p-1)$. Como aquí
5 de las 8 covariables no son significativas, el $R^2_{aj}$ es el que permite comparar honestamente con
modelos más pequeños (como el que sale del backward).

\res{$R^2=0,168$ y $R^2_{aj}=0,151$: el modelo tiene poder explicativo bajo.}

\pregunta{d)}{Supuestos del modelo}{Evalúe los principales supuestos del modelo lineal (linealidad,
normalidad de los residuos y homocedasticidad), incluyendo gráficos de diagnóstico y comentarios.}
\figcod{rlm-diag1}{Residuos vs.\ ajustados (arriba) y Q-Q normal (abajo).}
\figcod{rlm-diag2}{Escala-localización (arriba) y distribución de la nota final (abajo).}
\codigo{rlm-supuestos}
\salida{rlm-supuestos}
\begin{center}\small
\begin{tabular}{llcc}
\toprule
Supuesto & Prueba & $p$ & Conclusión \\ \midrule
Linealidad & RESET (potencias 2 y 3) & 0,020 & se rechaza al 5\,\% \\
Normalidad & Shapiro-Wilk & $1\times10^{-8}$ & se rechaza \\
Homocedasticidad & Breusch-Pagan & 0,020 & se rechaza al 5\,\% \\ \bottomrule
\end{tabular}
\end{center}
En el gráfico de residuos vs.\ ajustados aparece una \textbf{banda diagonal} de puntos con residuos de
$-5$ a $-13$: son los 38 ceros (residuo $=0-\hat y_i$). Esa banda inclina la curva LOESS (no linealidad),
produce la cola izquierda pesada del Q-Q y hace que la dispersión cambie con $\hat y$. Fuera de ella la
nube es bastante homogénea. Los estimadores MCO siguen siendo insesgados, pero los errores estándar y
los valores $p$ deben leerse con cautela.

\res{Ninguno de los tres supuestos se cumple al 5\,\%, y la causa principal es la masa de ceros en la
respuesta.}

\pregunta{e)}{Multicolinealidad (VIF)}{Analice la posible presencia de multicolinealidad utilizando el factor
de inflación de la varianza (VIF).}
\codigo{rlm-vif}
\salida{rlm-vif}
\[
\text{VIF}_j=\frac{1}{1-R_j^2}
\]
Todos los VIF están entre 1,07 y 1,74, muy por debajo de los umbrales usuales (5 o 10). Los más altos son
\var{educacion_madre} (1,74) y \var{educacion_padre} (1,70), que tienen correlación 0,62 entre sí; aun así
la varianza de sus coeficientes solo se infla un 70\,\%.

\res{No hay multicolinealidad preocupante. Que \var{educacion_padre} no sea significativa no se debe a
colinealidad, sino a que su efecto ya lo recoge \var{educacion_madre}.}

\pregunta{Backward a)}{Variables eliminadas y modelo final}{Aplique backward elimination con AIC partiendo
del modelo completo. Indique qué variables son eliminadas y presente el modelo final seleccionado.}
\codigo{rlm-backward}
\salida{rlm-backward}
\begin{center}\small
\begin{tabular}{clc}
\toprule
Paso & Variable eliminada & AIC tras eliminarla \\ \midrule
0 & (modelo completo) & 1146,5 \\
1 & educacion\_padre & 1144,6 \\
2 & internet\_casa & 1143,5 \\
3 & tiempo\_estudio & 1142,8 \\
4 & inasistencias & 1142,3 \\ \bottomrule
\end{tabular}
\end{center}
En el paso 5 quitar cualquier variable sube el AIC (edad: 1143, apoyo\_familiar: 1144), así que el algoritmo
se detiene. \var{step()} reporta el AIC sin las constantes de la verosimilitud (\var{extractAIC}); las
diferencias son las mismas que con \var{AIC()}.
\codigo{rlm-final}
\salida{rlm-final}
\[
\begin{aligned}\widehat{\text{nota\_final}}&=14,354-0,267\,\text{edad}-1,924\,\text{reprobaciones\_previas}\\
&\quad+0,621\,\text{educacion\_madre}-0,884\,\text{apoyo\_familiar}_{\text{sí}}\end{aligned}
\]

\res{Se eliminan educacion\_padre, internet\_casa, tiempo\_estudio e inasistencias. El modelo final conserva
edad, reprobaciones\_previas, educacion\_madre y apoyo\_familiar.}

\pregunta{Backward b)}{Comparación de AIC y parsimonia}{Compare el AIC del modelo completo con el del modelo reducido
y comente si la reducción implica una mejora en términos de parsimonia.}
\[
\text{AIC}=-2\ln L+2k
\]
\[
\text{AIC}_{\text{completo}}=2269,5\;(k=10),\qquad
\text{AIC}_{\text{reducido}}=2265,2\;(k=6),\qquad \Delta=4,3
\]
Con 4 parámetros menos el AIC baja 4,3 unidades: el modelo reducido pierde muy poco ajuste (la SCE pasa de
6877 a 6943) a cambio de ser más simple. El test $F$ parcial de las 4 variables eliminadas
($F_{4,386}=0,92$, $p=0,45$) confirma que en conjunto no aportan.

\res{Sí hay mejora en parsimonia: menor AIC con la mitad de covariables y sin pérdida significativa de ajuste.}

\pregunta{Backward c)}{$R^2$ ajustado del modelo final}{Analice el coeficiente de determinación ajustado del modelo
final y discuta los cambios respecto al modelo inicial.}
\begin{center}\small
\begin{tabular}{lccc}
\toprule
Modelo & covariables & $R^2$ & $R^2_{aj}$ \\ \midrule
Completo & 8 & 0,1684 & 0,1512 \\
Reducido (backward) & 4 & 0,1600 & 0,1519 \\ \bottomrule
\end{tabular}
\end{center}
El $R^2$ baja (0,168 $\to$ 0,160) porque siempre baja al quitar variables, pero el $R^2_{aj}$ sube
ligeramente (0,1512 $\to$ 0,1519): las variables eliminadas no compensaban los grados de libertad que
consumían. Además, \var{edad} y \var{apoyo_familiar} quedan con errores estándar algo menores.

\res{El $R^2_{aj}$ se mantiene (incluso sube un poco) con la mitad de variables; el modelo reducido explica lo
mismo con menos.}

\pregunta{Backward d)}{Supuestos del modelo seleccionado}{Evalúe los supuestos del modelo seleccionado e interprete
los resultados.}
\figcod{rlm-diag-final}{Diagnóstico del modelo reducido.}
\codigo{rlm-final-sup}
\salida{rlm-final-sup}
\begin{center}\small
\begin{tabular}{lccc}
\toprule
Prueba & Completo & Reducido & Conclusión (reducido) \\ \midrule
RESET (linealidad) & 0,020 & 0,030 & se rechaza al 5\,\% \\
Shapiro-Wilk (normalidad) & $1\times10^{-8}$ & $4\times10^{-9}$ & se rechaza \\
Breusch-Pagan (homocedasticidad) & 0,020 & 0,50 & \textbf{se cumple} \\
VIF máximo & 1,74 & 1,11 & sin multicolinealidad \\ \bottomrule
\end{tabular}
\end{center}
Al quitar variables desaparece la heterocedasticidad, pero la banda de ceros sigue ahí y mantiene la falta de
normalidad y el RESET significativo. Para comprobar que esa es la causa, reajusté el modelo final sin los 38
ceros:
\codigo{rlm-sensibilidad}
\salida{rlm-sensibilidad}
Sin los ceros los residuos pasan Shapiro-Wilk ($p=0,10$) y todos los signos se mantienen; el efecto de
\var{reprobaciones_previas} baja de $-1,92$ a $-1,21$ porque parte de él era ``reprobar $\Rightarrow$
abandonar''.

\res{El modelo reducido cumple homocedasticidad y no tiene multicolinealidad; la no normalidad se debe a los
alumnos que abandonaron. Concluyo que \textbf{reprobaciones previas} es el factor de mayor peso (negativo),
seguido de la \textbf{educación de la madre} (positivo); el apoyo familiar tiene signo negativo, que leo como
causalidad inversa (las familias apoyan más a quien ya va mal). Un modelo en dos partes (abandono sí/no + nota
si rinde) trataría mejor los ceros.}

\seccion{heartbeat}{Regresión logística binaria}{Regresión logística\\ binaria}{Heart Disease (Cleveland) $\cdot$ preguntas a -- g}
\section*{Datos y modelo}
\[
\log\frac{P(Y_i=1\mid x_i)}{1-P(Y_i=1\mid x_i)}=\beta_0+\beta_1\,\text{edad}_i+\beta_2\,\text{sexo}_i
+\beta_3\,\text{colesterol}_i+\beta_4\,\text{presion\_reposo}_i
\]
\[
+\,\beta_5\,\text{frecuencia\_cardiaca\_max}_i+\beta_6\,\text{angina\_ejercicio}_i+\beta_7\,\text{depresion\_st}_i
\]
con $Y_i=1$ si el paciente tiene enfermedad cardíaca. \var{sexo} (referencia: mujer) y
\var{angina_ejercicio} (referencia: no) entran como factores.
\codigo{log-datos}
\salida{log-datos}
Las clases están casi balanceadas: 137 enfermos (46,1\,\%) y 160 sanos.

\paragraph{Justificación de las variables.} Antes de ajustar comparé cada covariable entre enfermos y sanos:
\codigo{log-justifica}
\salida{log-justifica}
Seis de las siete variables difieren claramente entre grupos: los enfermos son mayores, con más presión,
menor frecuencia cardíaca máxima, mayor depresión ST, más hombres (56\,\% vs.\ 26\,\% de enfermos) y más
angina de esfuerzo (76\,\% vs.\ 32\,\%). El colesterol no difiere de forma marginal ($p=0,17$), pero lo
mantengo porque es un factor de riesgo clásico y puede aportar ajustado por las demás. Todas son medidas
clínicas baratas, disponibles antes de un cateterismo.

\pregunta{a)}{Ajuste por máxima verosimilitud}{Ajuste el modelo de regresión logística binaria mediante el
método de máxima verosimilitud.}
\[
\ln L(\boldsymbol\beta)=\sum_{i=1}^{n}\Big[y_i\,x_i^\top\boldsymbol\beta-\ln\!\big(1+e^{x_i^\top\boldsymbol\beta}\big)\Big],
\qquad
\frac{\partial\ln L}{\partial\boldsymbol\beta}=\sum_{i=1}^{n}x_i\,(y_i-\pi_i)=\mathbf 0
\]
No tiene solución cerrada: \var{glm()} la resuelve por Fisher scoring (IRLS), que convergió en 5 iteraciones.
\codigo{log-ajuste}
\salida{log-ajuste}

\res{La devianza baja de 409,95 (modelo nulo) a 274,16 con 7 parámetros adicionales; AIC $=290,2$.}

\pregunta{b)}{Tabla de estimación y significancia}{Presente la tabla de estimación de los parámetros e
identifique las variables estadísticamente significativas al nivel $\alpha=0,05$.}
\begin{center}\small
\begin{tabular}{lrrrrc}
\toprule
Variable & $\hat\beta$ & EE & $z$ & $p$ & ¿Signif.? \\ \midrule
Intercepto & $-2,383$ & 2,065 & $-1,15$ & 0,249 & no \\
edad & 0,0199 & 0,0195 & 1,02 & 0,309 & no \\
sexo (hombre) & 1,652 & 0,363 & 4,55 & $<0,001$ & \textbf{sí} \\
colesterol & 0,0062 & 0,0030 & 2,04 & 0,041 & \textbf{sí} \\
presion\_reposo & 0,0142 & 0,0091 & 1,56 & 0,118 & no \\
frecuencia\_cardiaca\_max & $-0,0296$ & 0,0084 & $-3,54$ & $<0,001$ & \textbf{sí} \\
angina\_ejercicio (sí) & 1,325 & 0,337 & 3,93 & $<0,001$ & \textbf{sí} \\
depresion\_st & 0,583 & 0,148 & 3,94 & $<0,001$ & \textbf{sí} \\ \bottomrule
\end{tabular}
\end{center}
Edad y presión en reposo sí diferían entre grupos por separado, pero dejan de ser significativas al
ajustar: su información ya la cargan la frecuencia cardíaca máxima (que baja con la edad) y la depresión ST.

\res{Significativas al 5\,\%: \textbf{sexo, colesterol, frecuencia\_cardiaca\_max, angina\_ejercicio y
depresion\_st}.}

\pregunta{c)}{Interpretación de edad, angina\_ejercicio y depresion\_st}{Interprete los coeficientes asociados
a las variables edad, angina\_ejercicio y depresion\_st en términos de odds ratios.}
\[
OR_j=e^{\hat\beta_j}=\frac{\text{odds}(x_j+1)}{\text{odds}(x_j)}
\]
\begin{itemize}[leftmargin=1.2em]
  \item[\color{cTeal}\faUser] \textbf{edad}: $e^{0,0199}=1,020$. Cada año adicional multiplica los odds de
  enfermedad por 1,02 (+2\,\%); en 10 años, $e^{10(0,0199)}=1,22$ (+22\,\%), con las demás variables fijas.
  No es significativo ($p=0,31$).
  \item[\color{cRed}\faHeartbeat] \textbf{angina\_ejercicio}: $e^{1,325}=3,76$. Un paciente con angina
  inducida por ejercicio tiene odds de enfermedad 3,8 veces mayores que uno sin angina, a igualdad del resto.
  \item[\color{cRed}\faChartLine] \textbf{depresion\_st}: $e^{0,583}=1,79$. Cada milímetro adicional de
  depresión del segmento ST multiplica los odds por 1,79 (+79\,\%).
\end{itemize}

\pregunta{d)}{Odds ratios e IC 95\,\%}{Calcule e interprete los odds ratios y sus intervalos de confianza al
95\,\%.}
\codigo{log-or}
\salida{log-or}
Los intervalos son por verosimilitud perfil (\var{confint} de un \var{glm}), mejores que $e^{\hat\beta\pm1,96\,EE}$
cuando la log-verosimilitud no es simétrica.
\begin{center}\small
\begin{tabular}{lrrl}
\toprule
Variable & OR & IC 95\,\% & Lectura \\ \midrule
sexo (hombre) & 5,22 & [2,62 ; 10,93] & odds 5 veces mayores en hombres \\
angina\_ejercicio (sí) & 3,76 & [1,96 ; 7,38] & odds casi 4 veces mayores \\
depresion\_st & 1,79 & [1,35 ; 2,42] & +79\,\% por mm \\
colesterol & 1,006 & [1,0003 ; 1,012] & +6,4\,\% cada 10 mg/dl \\
frecuencia\_cardiaca\_max & 0,971 & [0,955 ; 0,986] & $-26$\,\% cada 10 lpm ($0,971^{10}=0,74$) \\
presion\_reposo & 1,014 & [0,997 ; 1,033] & incluye 1: no concluyente \\
edad & 1,020 & [0,982 ; 1,060] & incluye 1: no concluyente \\ \bottomrule
\end{tabular}
\end{center}

\res{Los IC que no contienen al 1 coinciden con las variables significativas. Los factores de riesgo más
fuertes son ser hombre y tener angina de esfuerzo; una frecuencia cardíaca máxima alta es protectora.}

\pregunta{e)}{Calidad del ajuste: razón de verosimilitudes y matriz de confusión}{Evalúe la calidad del ajuste
utilizando la razón de verosimilitudes, la matriz de confusión y la tasa de clasificación correcta (umbral
0,5); interprete exactitud, sensibilidad y especificidad en el contexto del diagnóstico.}
\codigo{log-lr}
\salida{log-lr}
\[
G^2=-2\big[\ln L_0-\ln L_1\big]=409,95-274,16=135,8\;\sim\;\chi^2_7,\qquad p<2\times10^{-16}
\]
Rechazo $H_0:\beta_1=\dots=\beta_7=0$: el modelo explica significativamente más que solo la prevalencia.
El $R^2$ de McFadden (0,33) está en el rango que se considera un buen ajuste (0,2--0,4).
\codigo{log-confusion}
\salida{log-confusion}
\begin{center}\small
\begin{tabular}{lcc}
\toprule
& \multicolumn{2}{c}{Clase predicha} \\ \cmidrule(lr){2-3}
Clase real & No & Sí \\ \midrule
No & VN $=133$ & FP $=27$ \\
Sí & FN $=37$ & VP $=100$ \\ \bottomrule
\end{tabular}
\end{center}
\[
\text{Exactitud}=\frac{100+133}{297}=0,785
\]
\[
\text{Sensibilidad}=\frac{100}{100+37}=0,730,\qquad
\text{Especificidad}=\frac{133}{133+27}=0,831
\]
\begin{itemize}[leftmargin=1.2em]
  \item[\color{cTeal}\faCheck] \textbf{Exactitud 78,5\,\%}: clasifica bien a casi 4 de cada 5 pacientes
  (la regla trivial ``todos sanos'' acertaría 53,9\,\%).
  \item[\color{cRed}\faExclamationTriangle] \textbf{Sensibilidad 73,0\,\%}: detecta 100 de 137 enfermos;
  \textbf{37 enfermos quedan como sanos}. En diagnóstico cardíaco este es el error caro, porque un enfermo
  no detectado no recibe tratamiento.
  \item[\color{cTeal}\faIcon{shield-alt}] \textbf{Especificidad 83,1\,\%}: 27 sanos serían derivados a pruebas
  adicionales sin necesidad; es un costo, pero menor.
\end{itemize}

\res{El modelo ajusta significativamente ($G^2=135,8$) y clasifica bien el 78,5\,\%. Con umbral 0,5 es más
específico que sensible; para tamizaje convendría bajar el umbral (ver pregunta f).}

\pregunta{f)}{Curva ROC y AUC}{Analice la capacidad predictiva del modelo mediante la curva ROC e interprete el
área bajo la curva (AUC).}
\figcod{log-roc}{Curva ROC con el punto de Youden (arriba) y probabilidades estimadas por clase (abajo).}
\codigo{log-roc}
\salida{log-roc}
$\text{AUC}=0,858$ (IC 95\,\% DeLong: 0,816--0,900): si tomo al azar un enfermo y un sano, en el 85,8\,\% de
los casos el modelo le asigna mayor probabilidad al enfermo. Es una discriminación \emph{buena} (0,8--0,9) y el
IC está lejos de 0,5. El umbral de Youden (0,443) sube la sensibilidad a 78,1\,\% con especificidad de 80\,\%,
un mejor equilibrio para tamizaje que 0,5.

\res{AUC $=0,858$: el modelo discrimina bien entre enfermos y sanos.}

\pregunta{g)}{Limitaciones y extensiones}{Discuta las limitaciones del modelo y proponga posibles extensiones.}
Todas las métricas anteriores se calcularon sobre los mismos 297 pacientes del ajuste, lo que las hace
optimistas. Lo verifiqué con validación cruzada de 10 pliegues:
\codigo{log-cv}
\salida{log-cv}
Fuera de muestra el AUC baja de 0,858 a 0,836 y la exactitud de 78,5\,\% a 76,1\,\%: hay algo de
sobreajuste, pero el modelo generaliza razonablemente.

\noindent\textbf{Limitaciones}
\begin{itemize}[leftmargin=1.2em]
  \item[\color{cRed}\faMinusCircle] Solo 297 pacientes de un único hospital (Cleveland): poca validez externa.
  \item[\color{cRed}\faMinusCircle] Supone que el logit es lineal en cada variable continua (edad, colesterol,
  frecuencia cardíaca, depresión ST) y no incluye interacciones (por ejemplo, sexo $\times$ edad).
  \item[\color{cRed}\faMinusCircle] Deja fuera predictores muy informativos de la base: tipo de dolor torácico,
  vasos coloreados por fluoroscopía, talasemia y pendiente ST.
  \item[\color{cRed}\faMinusCircle] Es un estudio transversal: los OR son asociaciones, no efectos causales.
\end{itemize}
\noindent\textbf{Extensiones}
\begin{itemize}[leftmargin=1.2em]
  \item[\color{cTeal}\faPlusCircle] Agregar las variables omitidas y seleccionar con AIC/BIC o con
  regularización (lasso).
  \item[\color{cTeal}\faPlusCircle] Términos no lineales (splines en edad y frecuencia cardíaca) o un GAM.
  \item[\color{cTeal}\faPlusCircle] Usar \var{diagnostico_original} (0--4) en un modelo logístico ordinal de
  severidad.
  \item[\color{cTeal}\faPlusCircle] Elegir el umbral según el costo relativo FN/FP y validar en una cohorte externa.
\end{itemize}

\seccion{project-diagram}{Análisis de componentes principales}{Análisis de\\ componentes principales}{Student Performance $\cdot$ preguntas a -- g}
\section*{Variables}
Uso las 11 variables cuantitativas del enunciado, todas como numéricas:
\codigo{acp-datos}
\salida{acp-datos}

\pregunta{a)}{Por qué estandarizar}{Explique por qué es necesario estandarizar las variables antes de aplicar el
ACP en este contexto.}
El ACP busca las direcciones de máxima varianza, así que una variable con varianza grande domina las
primeras componentes solo por su escala. Aquí \var{inasistencias} va de 0 a 75 con $s=8,0$ ($s^2=64$),
mientras que las escalas Likert 1--5 tienen $s\approx1$ y \var{reprobaciones_previas} $s=0,74$ ($s^2=0,55$).
Sin estandarizar, la CP1 sería prácticamente ``inasistencias'' (64 de unas 76 unidades de varianza
total, el 84\,\%). Además las unidades no son comparables (años, faltas, escalas ordinales). Con
\[
z_{ij}=\frac{x_{ij}-\bar x_j}{s_j}
\]
todas tienen varianza 1 y el ACP sobre $Z$ equivale a diagonalizar la matriz de correlaciones $R$.

\res{Estandarizo para que cada variable pese lo mismo y las componentes reflejen correlaciones, no unidades.}

\pregunta{b)}{Matriz de correlaciones}{Calcule la matriz de correlaciones entre las variables y comente
brevemente la presencia de relaciones lineales relevantes.}
\codigo{acp-cor}
\salida{acp-cor}
\figcod{acp-corplot}{Matriz de correlaciones (triángulo superior).}
Relaciones más relevantes:
\begin{itemize}[leftmargin=1.2em]
  \item[\color{cTeal}\faLink] \textbf{alcohol días de semana -- fin de semana}: $r=0,65$, la más fuerte.
  \item[\color{cTeal}\faLink] \textbf{educación de la madre -- del padre}: $r=0,62$.
  \item[\color{cTeal}\faLink] \textbf{salidas sociales} con alcohol fin de semana ($0,42$), tiempo libre ($0,29$)
  y alcohol semana ($0,27$): un bloque de ``vida social''.
  \item[\color{cRed}\faUnlink] \textbf{reprobaciones} con educación de los padres ($-0,24$, $-0,25$) y con
  edad ($+0,24$); \textbf{tiempo de estudio} con alcohol fin de semana ($-0,25$).
\end{itemize}
El resto está por debajo de $|0,2|$ y \var{salud} casi no se relaciona con nada. El test de esfericidad de
Bartlett ($\chi^2_{55}=701,3$, $p\approx10^{-113}$) rechaza $R=I$: hay correlación suficiente para que el ACP
tenga sentido, aunque las correlaciones moderadas anticipan que la varianza se repartirá en varias componentes.

\pregunta{c)}{ACP con la matriz de correlaciones}{Aplique el Análisis de Componentes Principales utilizando la
matriz de correlaciones.}
\[
R\,\mathbf a_k=\lambda_k\,\mathbf a_k,\qquad \sum_{k=1}^{11}\lambda_k=\operatorname{tr}(R)=11,\qquad
\text{CP}_k=Z\,\mathbf a_k
\]
\codigo{acp-pca}
\salida{acp-pca}
Con \var{prcomp(scale. = TRUE)} los autovalores son los de $R$ (suman 11) y los vectores \var{rotation}
son las cargas $\mathbf a_k$.

\pregunta{d)}{Scree plot y número de componentes}{Presente el scree plot e indique cuántos componentes deberían
retenerse según la varianza explicada acumulada (80\,\% o 90\,\%).}
\figcod{acp-scree}{Scree plot (arriba) y varianza acumulada con las líneas de 80\,\% y 90\,\% (abajo).}
\begin{center}\small
\begin{tabular}{lcc}
\toprule
Criterio & Componentes & Varianza retenida \\ \midrule
Varianza acumulada $\ge$ 80\,\% & \textbf{7} & 81,9\,\% \\
Varianza acumulada $\ge$ 90\,\% & 9 & 94,0\,\% \\
Kaiser ($\lambda_k>1$) & 4 & 58,3\,\% \\
Codo del scree plot & 3 & 48,8\,\% \\ \bottomrule
\end{tabular}
\end{center}
Retengo \textbf{7 componentes} (criterio del 80\,\%): con 9 llegaría al 94\,\%, pero pasaría de 11 a 9
variables, que casi no reduce la dimensión. La curva cae despacio después de CP2, consecuencia de las
correlaciones bajas: este conjunto no se comprime bien en pocas componentes.

\res{7 CP según el 80\,\% (81,9\,\%); 9 según el 90\,\%.}

\pregunta{e)}{Interpretación de CP1 y CP2}{Interprete los dos primeros componentes en términos de las variables
originales, analizando los signos y magnitudes de las cargas.}
\codigo{acp-cargas}
\salida{acp-cargas}
\figcod{acp-cargas}{Cargas de las dos primeras componentes.}
\begin{itemize}[leftmargin=1.2em]
  \item[\color{cTeal}\faGlassCheers] \textbf{CP1 (21,1\,\%) -- ``estilo de vida social / riesgo''}: cargas
  negativas grandes en alcohol fin de semana ($-0,52$), alcohol semana ($-0,48$), salidas sociales ($-0,37$),
  reprobaciones ($-0,30$), tiempo libre ($-0,27$) y edad ($-0,24$), y positiva en tiempo de estudio ($+0,28$).
  Un CP1 \emph{bajo} es un alumno que bebe, sale mucho, estudia poco, es mayor y ha reprobado; un CP1
  \emph{alto} es el perfil estudioso. (El signo global es arbitrario.)
  \item[\color{cTeal}\faUserGraduate] \textbf{CP2 (16,6\,\%) -- ``capital educativo familiar''}: dominada por
  educación de la madre ($+0,61$) y del padre ($+0,60$), contra reprobaciones ($-0,30$) y edad ($-0,22$). Un
  CP2 alto es un alumno de padres con educación superior y sin repitencia.
\end{itemize}
\var{salud} e \var{inasistencias} casi no cargan en ninguna de las dos; aparecen en CP3 (inasistencias
$-0,61$, salud $+0,44$).

\pregunta{f)}{Individuos en el plano CP1--CP2}{Represente gráficamente los individuos en el plano definido por
los dos primeros componentes e identifique posibles patrones o agrupamientos.}
\figcod{acp-individuos}{Puntajes de los 395 alumnos en CP1--CP2.}
\figcod{acp-biplot}{Biplot: individuos (gris) y cargas (flechas).}
\begin{itemize}[leftmargin=1.2em]
  \item[\color{cTeal}\faEye] No hay grupos separados: la nube es continua y asimétrica, con una cola hacia
  CP1 negativo.
  \item[\color{cRed}\faEye] Esa cola son los alumnos con alto consumo de alcohol el fin de semana (rojo): su
  CP1 medio es $-1,98$ frente a $0,50$ del resto. El biplot muestra las flechas de alcohol y salidas
  apuntando justo hacia ella.
  \item[\color{cTeal}\faEye] Por aprobación los puntos se mezclan, pero los aprobados están algo desplazados
  hacia arriba y a la derecha (CP1 medio $0,23$ vs.\ $-0,47$; CP2 $0,15$ vs.\ $-0,30$): estudiar más, beber
  menos y tener padres más educados va con aprobar.
\end{itemize}

\pregunta{g)}{CP como variables explicativas}{Discuta si los componentes principales obtenidos podrían ser
utilizados posteriormente como variables explicativas en un modelo de regresión.}
Sí: es la \emph{regresión sobre componentes principales}. Lo probé con las 7 CP retenidas para explicar la nota
final:
\codigo{acp-regresion}
\salida{acp-regresion}
\begin{itemize}[leftmargin=1.2em]
  \item[\color{cTeal}\faPlusCircle] \textbf{Ventaja}: las CP son ortogonales (VIF $=1$ exacto), así que
  eliminan cualquier multicolinealidad y cada coeficiente se estima de forma independiente.
  \item[\color{cTeal}\faPlusCircle] CP1 y CP2 resultan significativas y con signo esperado: más ``estilo
  estudioso'' y más capital educativo, mejor nota.
  \item[\color{cRed}\faMinusCircle] \textbf{Desventaja 1}: las CP se eligen por varianza de las $X$, no por
  su relación con $Y$. Aquí CP5 y CP7, de poca varianza, son significativas, y CP3--CP4 no aportan.
  \item[\color{cRed}\faMinusCircle] \textbf{Desventaja 2}: interpretar ``+1 en CP7'' es difícil, y con 7 de
  11 dimensiones se pierde algo de ajuste ($R^2_{aj}$: 0,137 frente a 0,153 con las 11 variables originales).
  \item[\color{cRed}\faMinusCircle] En esta base la multicolinealidad no era un problema (en el ejercicio de regresión lineal
  sobre esta base los VIF fueron $<1,8$), por lo que el beneficio principal del ACP no aplica.
\end{itemize}

\res{Se pueden usar y garantizan ortogonalidad, pero aquí convienen poco: las correlaciones son bajas, se pierde
interpretabilidad y ajuste. Las CP serían útiles con muchas variables muy correlacionadas o con $p$ cercano a
$n$; una alternativa es PLS, que elige componentes pensando en $Y$.}

\seccion{object-group}{Análisis de conglomerados}{Análisis de\\ conglomerados}{Heart Disease (Cleveland) $\cdot$ Ward.D2, $k$-means y comparación}
\section*{Variables}
Uso \var{edad}, \var{presion_reposo}, \var{colesterol}, \var{frecuencia_cardiaca_max} y
\var{depresion_st}, estandarizadas.
\codigo{clu-datos}
\salida{clu-datos}

\pregunta{a)}{Por qué estandarizar}{Explique por qué es necesario estandarizar las variables antes de realizar
un análisis de conglomerados en este contexto.}
La distancia euclidiana $d(x_i,x_j)=\sqrt{\sum_k (x_{ik}-x_{jk})^2}$ suma diferencias en unidades
originales. El colesterol tiene $s=52$ mg/dl y la depresión ST $s=1,17$ mm: una diferencia típica de
colesterol pesa unas $(52/1,17)^2\approx 2000$ veces más que una de ST en $d^2$. Sin estandarizar, los grupos
serían básicamente ``colesterol alto / bajo''. Lo comprobé: la correlación entre la distancia sin estandarizar
y la diferencia de colesterol es 0,94 (figura de la pregunta b).

\res{Estandarizo ($z=(x-\bar x)/s$) para que las cinco variables clínicas contribuyan por igual a la distancia.}

\pregunta{b)}{Matriz de distancias euclidianas}{Calcule la matriz de distancias euclidianas entre los individuos
a partir de las variables estandarizadas.}
\codigo{clu-dist}
\salida{clu-dist}
La matriz es $297\times297$ simétrica con diagonal cero; tiene $\binom{297}{2}=43\,956$ distancias distintas.
Muestro los 6 primeros pacientes: el 2 y el 4 son los más distintos ($d=5,40$) y el 1 y el 3 los más
parecidos ($d=1,76$). La distancia mediana es 2,88.
\figcod{clu-dist}{Distribución de distancias (arriba) y distancia estandarizada vs.\ sin estandarizar (abajo).}

\pregunta{c)}{Ward.D2}{Aplique un análisis de conglomerados jerárquicos utilizando el método de Ward (ward.D2).}
En cada paso Ward une los dos grupos que menos aumentan la suma de cuadrados dentro de grupos:
\[
\Delta(A,B)=\frac{n_An_B}{n_A+n_B}\,\lVert\bar z_A-\bar z_B\rVert^2
\]
\var{ward.D2} trabaja con las distancias euclidianas sin elevar (las eleva internamente), que es la
implementación correcta del criterio de Ward.
\codigo{clu-ward}
\salida{clu-ward}

\pregunta{d)}{Dendrograma y número de grupos}{Represente el dendrograma e indique el número de conglomerados
que considera adecuado, justificando su elección.}
\figcod{clu-dendro}{Dendrograma Ward.D2 (arriba) y altura de las últimas fusiones (abajo).}
Elijo \textbf{$k=3$}:
\begin{itemize}[leftmargin=1.2em]
  \item[\color{cTeal}\faCut] Las dos últimas fusiones son los saltos grandes: 23,4 (de 2 a 1 grupo) y 18,2 (de
  3 a 2). Desde $k=4$ las fusiones bajan en pasos de apenas $\approx1,3$: el codo está en $k=3$.
  \item[\color{cTeal}\faCut] En el dendrograma, un corte a altura $\approx 16$ cruza tres ramas largas y de
  tamaño parecido (107, 100 y 90 pacientes).
  \item[\color{cRed}\faCut] La silueta media favorece algo a $k=2$ (0,20 vs.\ 0,15). Aun así prefiero $k=3$
  porque separa un perfil clínico más (el de presión y colesterol altos, ver pregunta e), que con $k=2$ se pierde.
\end{itemize}

\pregunta{e)}{Asignación y características de cada grupo}{Asigne cada individuo a un conglomerado y describa
las principales características clínicas de cada grupo.}
\codigo{clu-perfil-hc}
\salida{clu-perfil-hc}
\begin{center}\small
\begin{tabular}{lccc}
\toprule
& \textbf{Grupo 1} & \textbf{Grupo 2} & \textbf{Grupo 3} \\ \midrule
$n$ & 107 & 100 & 90 \\
Edad (años) & 60,3 & 48,0 & 55,0 \\
Presión reposo (mmHg) & 132,0 & 124,9 & \textbf{138,9} \\
Colesterol (mg/dl) & 239,7 & 215,4 & \textbf{291,9} \\
FC máxima (lpm) & \textbf{131,5} & 163,1 & 156,0 \\
Depresión ST (mm) & \textbf{2,0} & 0,5 & 0,6 \\
\% hombres & 76,6 & 72,0 & 52,2 \\
\% angina de esfuerzo & 51,4 & 16,0 & 28,9 \\ \midrule
\% con enfermedad & \textbf{72,9} & 24,0 & 38,9 \\ \bottomrule
\end{tabular}
\end{center}
\begin{itemize}[leftmargin=1.2em]
  \item[\color{cRed}\faHeartbeat] \textbf{Grupo 1 -- isquémico}: los mayores, con baja frecuencia cardíaca
  máxima y depresión ST alta (signos de isquemia en la prueba de esfuerzo); la mitad tiene angina.
  \item[\color{cTeal}\faRunning] \textbf{Grupo 2 -- joven y sano}: 48 años en promedio, buena capacidad
  cardíaca (163 lpm), ST casi normal y los valores más bajos de presión y colesterol.
  \item[\color{cPurple}\faTint] \textbf{Grupo 3 -- metabólico}: colesterol muy alto (292) y presión alta, pero
  con prueba de esfuerzo casi normal; tiene la mayor proporción de mujeres.
\end{itemize}

\pregunta{f)}{$k$-means con $k=3$}{Aplique el método de $k$-means considerando el mismo número de conglomerados
seleccionado en el método jerárquico.}
$k$-means minimiza $\sum_{g=1}^{3}\sum_{i\in g}\lVert z_i-\bar z_g\rVert^2$. Uso 50 arranques aleatorios
(\var{nstart = 50}) para no quedarme en un mínimo local, y reetiqueto los grupos para que el grupo $j$ sea el
que más se parece al grupo $j$ de Ward.
\codigo{clu-kmeans}
\salida{clu-kmeans}
Los tres grupos explican el 35,3\,\% de la variabilidad total ($SC_{entre}/SC_{total}$).

\pregunta{g)}{Centroides y perfiles de riesgo}{Analice los centroides obtenidos e interprete los perfiles de cada
grupo en términos de riesgo cardiovascular.}
\figcod{clu-centroides}{Centroides de $k$-means en unidades $z$.}
\begin{center}\small
\begin{tabular}{lccccc|c}
\toprule
Grupo & Edad & Presión & Colest. & FC máx. & ST & \% enfermos \\ \midrule
1 ($n=85$) & 59,2 & 127,1 & 228,6 & \textbf{125,9} & \textbf{1,9} & \textbf{74,1} \\
2 ($n=142$) & 48,8 & 124,9 & 236,4 & 164,3 & 0,4 & 23,2 \\
3 ($n=70$) & 60,5 & \textbf{151,0} & \textbf{292,3} & 148,6 & 1,3 & 58,6 \\ \bottomrule
\end{tabular}
\end{center}
\begin{itemize}[leftmargin=1.2em]
  \item[\color{cRed}\faExclamationTriangle] \textbf{Grupo 1 -- riesgo alto (isquémico)}: centroide con FC
  máxima a $-1,03$ desviaciones y ST a $+0,72$. Es la respuesta anormal al esfuerzo típica de enfermedad
  coronaria; 3 de cada 4 están enfermos.
  \item[\color{cTeal}\faCheckCircle] \textbf{Grupo 2 -- riesgo bajo}: todos los valores por debajo de la media
  salvo la FC máxima ($+0,64$). Jóvenes con buena reserva cardíaca; 23\,\% de enfermos.
  \item[\color{cPurple}\faExclamationCircle] \textbf{Grupo 3 -- riesgo intermedio-alto (factores clásicos)}:
  presión ($+1,09$) y colesterol ($+0,87$) altos y mayor edad. Aún sin daño evidente en la prueba de esfuerzo,
  el 59\,\% ya está enfermo: es el grupo de prevención.
\end{itemize}

\pregunta{Comparación a)}{Ward.D2 vs.\ $k$-means}{Compare los conglomerados obtenidos mediante los métodos
ward.D2 y $k$-means.}
\codigo{clu-compara}
\salida{clu-compara}
\figcod{clu-plano}{Los dos agrupamientos en el plano de las dos primeras CP de las 5 variables.}
\begin{itemize}[leftmargin=1.2em]
  \item[\color{cTeal}\faCheck] Coinciden en el 73,4\,\% de los pacientes (78 + 91 + 49 de 297). El índice de
  Rand ajustado es 0,40: acuerdo moderado.
  \item[\color{cTeal}\faCheck] Los grupos 1 y 2 son casi los mismos (73\,\% y 91\,\% de coincidencia).
  \item[\color{cRed}\faIcon{exchange-alt}] La diferencia está en el grupo 3: Ward lo forma con todo el colesterol
  alto (90 pacientes); $k$-means pasa 41 de ellos al grupo 2 y se queda con
  un grupo 3 más extremo (presión 151, colesterol 292). Por eso $k$-means queda con un grupo 2 más grande (142).
  \item[\color{cTeal}\faBalanceScale] $k$-means logra grupos más compactos: silueta 0,23 frente a 0,15 de Ward.
\end{itemize}

\pregunta{Comparación b)}{Similitudes y diferencias entre enfoques}{Discuta las similitudes y diferencias entre ambos
enfoques.}
\begin{center}\small
\begin{tabular}{p{3.3cm}p{5.4cm}p{5.4cm}}
\toprule
& \textbf{Ward.D2 (jerárquico)} & \textbf{$k$-means (no jerárquico)} \\ \midrule
Criterio & Minimiza el aumento de SC dentro en cada fusión & Minimiza la SC dentro total \\
Número de grupos & Se decide después, mirando el dendrograma & Hay que fijarlo antes \\
Asignaciones & Definitivas: una fusión no se deshace & Se reasignan en cada iteración \\
Resultado & Determinista & Depende del arranque (uso 50) \\
Costo & $O(n^2)$ en memoria (matriz de distancias) & Lineal en $n$, escala a bases grandes \\
Aquí & Silueta 0,15; grupos 107/100/90 & Silueta 0,23; grupos 85/142/70 \\ \bottomrule
\end{tabular}
\end{center}
Ambos comparten la distancia euclidiana, las variables estandarizadas y la misma función objetivo (suma de
cuadrados dentro), por eso llegan a perfiles parecidos. En la práctica son complementarios: Ward sirve para
elegir $k$ y $k$-means para refinar las asignaciones.

\pregunta{Comparación c)}{Relación con enfermedad\_cardiaca}{Analice la relación entre los conglomerados obtenidos y la
variable enfermedad\_cardiaca.}
\codigo{clu-enfermedad}
\salida{clu-enfermedad}
\figcod{clu-enfermedad}{Porcentaje de enfermos por grupo (línea roja: prevalencia global, 46,1\,\%).}
\begin{itemize}[leftmargin=1.2em]
  \item[\color{cTeal}\faCheck] La relación es fuerte en ambos métodos: $\chi^2_2=52$ (Ward) y $\chi^2_2=61$
  ($k$-means), con $p<10^{-11}$. La V de Cramer de $k$-means es 0,45 (asociación moderada a fuerte).
  \item[\color{cTeal}\faCheck] El grupo 1 tiene 73--74\,\% de enfermos y el grupo 2, 23--24\,\%: una diferencia
  de 3 veces, sin haberle mostrado nunca la etiqueta al algoritmo.
  \item[\color{cRed}\faIcon{exchange-alt}] En el grupo 3 los métodos difieren (39\,\% vs.\ 59\,\%): el grupo 3 más
  extremo de $k$-means concentra más enfermos, por eso su asociación es más fuerte.
  \item[\color{cTeal}\faBalanceScale] Comparado con la regresión logística sobre esta misma base (AUC 0,86), el enfoque no supervisado separa
  menos, porque no usa sexo ni angina y no optimiza la predicción; a cambio entrega perfiles clínicos
  interpretables (isquémico, sano, metabólico).
\end{itemize}

\res{Los conglomerados, construidos solo con 5 variables clínicas, recuperan buena parte de la estructura de
la enfermedad: el perfil ``isquémico'' (baja FC máxima, ST alta) concentra a los enfermos y el perfil
``joven sano'' a los sanos.}

\end{document}
